\begin{abstract}
Mobile applications increasingly mediate essential financial transactions and citizen services, exposing them to stringent statutory data-protection and cybersecurity mandates. However, automated compliance tools frequently treat statutory verification as generic vulnerability scanning, overlooking the fundamental boundary between observable client software behavior and unobservable backend or organizational processes. In this paper, we present \textbf{Reg2App}, a formally specified, provenance-preserving methodology for translating natural-language regulatory requirements into partially observable Android program properties. Reg2App decouples legal obligations from technical controls through a 7-tuple mapping model, bounds analysis over a two-dimensional observability taxonomy, and evaluates conformance over a discrete 5-state epistemic assessment space that preserves evidence provenance. Crucially, Reg2App enforces a principled three-tier distinction between primary statutory mandates (e.g., transport encryption), operationalized supervisory baseline controls (e.g., certificate pinning), and empirical security-risk indicators (e.g., foreign telemetry endpoints). We instantiate our approach using the emerging regulatory corpus of Bangladesh (central bank directives, data protection legislation, and cybersecurity statutes). A preliminary expert-elicitation pilot ($m=3$ raters, 12 mappings, pre-reconciliation raw agreement $88.9\%$, $\text{Fleiss' } \kappa = 0.776$, $\text{Krippendorff's } \alpha = 0.782$, reaching 100\% unanimous consensus via structured Delphi reconciliation) demonstrates the operational feasibility of formalizing regulatory clauses. We verify analyzer implementation correctness using a 30-application synthetic micro-benchmark regression suite. On real-world production bytecode, independent manual adjudication with certified external auditors confirmed $96.0\%$ precision in a finding-sample study of 50 tool-reported verdicts ($48/50$, $95\%$ Wilson CI: $[86.5\%, 98.9\%]$) and $94.1\%$ recall in an independent audit of 51 ground-truth controls across five applications ($48/51$, $95\%$ Wilson CI: $[83.8\%, 97.9\%]$, defect recall $85.0\%$). We then conduct an empirical ecosystem study across 53 production Android applications sampled in early 2026 across eight economic sectors in Bangladesh, comprising 39 live financial applications (all 14 operational Mobile Financial Services wallets and 25 commercial banks) contrasted observationally against 14 non-financial consumer applications (revealing an observational baseline disparity—not causal regulatory effect—where retail e-commerce exhibited zero observed conformance across 6 applicable rules, Mann-Whitney $U = 156.0, p < 0.001$, reflecting differing supervisory scopes and resources). Our static analysis demonstrates that while $100\%$ of financial apps enforce statutory transport encryption (TLS), $82.1\%$ lack supervisory defense-in-depth certificate pinning, and $30.8\%$ leave backup extraction flags enabled—evaluated under strict offline ethical auditing and institutional pseudonymization. Finally, a counterbalanced laboratory experiment with 60 developers and auditors demonstrates that evidence-grounded bilingual reporting improves comprehension ($53.6\% \rightarrow 82.9\%$) and accelerates defect remediation by $52.1\%$ (from $39.7$ to $19.0$ minutes, $d = 2.45$) over unparsed scanner traces. All Regulation-as-Code ontologies, benchmarks, and pipelines are released for open reproduction.
\end{abstract}

\begin{keyword}
Software Engineering \sep Android Security \sep Regulatory Compliance \sep Regulation-as-Code \sep Evidence-Based Verification \sep Epistemic Assessment \sep Program Analysis
\end{keyword}
